\emph{(i)} On a generator the boost of every horizon in the scope of Premise \claimref{synthesis.n-plus-null-ground-state} acts as $u \mapsto e^{\kappa t}u$ (the setup of Theorem \claimref{synthesis.positive-null-energy-gives-r-premises}, the boost $\beta_{t}$: the Killing field is $\kappa u\,\partial_{u}$ on the generator for the eternal hole and for de Sitter space, and the local boost rotor dilates each generator of the local Rindler plane about its edge point). On fields it is $\beta_{t}(\psi(u)) = e^{\kappa t/2}\psi(e^{\kappa t}u)$, which is $\sigma^{0}_{\lambda}$ at $\lambda = e^{\kappa t}$. For $\sigma_{\lambda}f(u) = \lambda^{p}f(\lambda u)$ one has $\ip{\sigma_{\lambda}f}{\sigma_{\lambda}g} = \lambda^{2p - 1}\ip{f}{g}$, so the anticommutator $\{a(f), a^{*}(g)\} = \ip{f}{g}$ is preserved exactly for $p = \tfrac12$; the check verifies this on indicator functions of rational intervals at $\lambda = (3/2)^{4}, (1/5)^{4}$ among the weights $p = k/4$, $0 \le k \le 4$. In $\Cl(1,3)$ with $Q(e_{0}) = 1$, $Q(e_{j}) = -1$, $(e_{1}e_{0})^{2} = 1$ and $B = \cosh\tfrac{\lambda}{2} + \sinh\tfrac{\lambda}{2}\,e_{1}e_{0}$ is a rotor; $e_{1}e_{0}$ anticommutes with $e_{0}, e_{1}$ and commutes with $e_{2}, e_{3}$, and $e_{1}e_{0}\ell = \ell$, $e_{1}e_{0}n = -n$, so $B\ell\tilde B = B^{2}\ell = e^{\lambda}\ell$, $Bn\tilde B = e^{-\lambda}n$, and $u\ell + y \mapsto e^{\lambda}u\ell + y$ for $y$ in the span of $e_{2}, e_{3}$; the check verifies these at $e^{\lambda/2} = \tfrac32, \tfrac15, 7$. So the boost is the scaling of the generator about the edge point with the weight $\tfrac12$ of a chiral field of scaling dimension $\tfrac12$.

\emph{(ii)} If $\omega\circ\beta_{t} = \omega$ for all $t$, then by (i) $\omega\circ\sigma^{0}_{\lambda} = \omega$ for all $\lambda > 0$, that is $\lambda W_{\omega}(\lambda x, \lambda y) = W_{\omega}(x, y)$; the family whose limit defines the scaling limit at the edge is constant, so the limit exists and is $\omega$. If $\omega$ is locally the null vacuum at the edge, then $W_{\omega} = W_{0}$, and a gauge-invariant quasi-free state is determined by its two-point function, so $\omega = \omega_{0}$. Conversely $W_{0}$ is invariant under translations and under $\sigma^{u_{0}}_{\lambda}$ for every $u_{0}$, so $\omega_{0}$ is its own scaling limit at every point. By Theorem \claimref{synthesis.positive-null-energy-gives-r-premises}(i) the gauge-invariant quasi-free states satisfying (N$_{+}$) are exactly $\omega_{0}$, which gives the stated equivalence. For the weaker form: $W_{0}$ is stationary, so $D = W_{\omega} - W_{0}$ satisfies $\lambda D(\lambda x, \lambda y) = D(x, y)$. If $|D| \le C$ on a neighbourhood $U$ of $(0, 0)$, then for any $(x, y)$ and $\lambda$ small enough that $(\lambda x, \lambda y) \in U$, $|D(x, y)| = \lambda|D(\lambda x, \lambda y)| \le \lambda C$, and $\lambda \to 0$ gives $D(x, y) = 0$.

\emph{(iii)} The two halves $\CAR(L^{2}(\R_{\pm}, du)\otimes K)$ generate $\CAR(L^{2}(\R, du)\otimes K)$ as a graded tensor product, and the product of two even states is an even state; each half is $\beta$-invariant, so the product of stationary states is stationary. On $u > 0$ put $u = e^{x}$ and $\phi(x) = e^{x/2}\psi(e^{x})$: the map is unitary from $L^{2}(\R_{+}, du)$ to $L^{2}(\R, dx)$, the boost becomes the translation $x \mapsto x + \kappa t$, and the null vacuum's restriction becomes $\tfrac{i}{2\pi}e^{(x + x')/2}(e^{x} - e^{x'} + i0)^{-1} = \tfrac{i}{2\pi}\big(2\sinh((x - x' + i0)/2)\big)^{-1}$, the chiral thermal two-point function at period $2\pi$ (Theorem \claimref{synthesis.positive-null-energy-gives-r-premises}(ii)). The chiral two-point function KMS at period $2\pi c$ for the translations is the same function of $(x - x')/c$ divided by $c$ (the leading singularity $\tfrac{i}{2\pi}(x - x' + i0)^{-1}$, fixed by the anticommutator, is then the same for every $c$), which gives $W_{c}$; its Fourier transform in $x - x'$ is that at $c = 1$ evaluated at $c\,w$, so the Mellin symbol is $m_{1}(c w) = (1 + e^{-2\pi c w})^{-1}$, the symbol of the I-state of Theorem \claimref{hawking.i-state-continuum} for $h = -\kappa w$ at $\beta = 2\pi c/\kappa$, which is jointly KMS by Fact \claimref{hawking.imported-quasi-free-kms-reduces-two}. In the check the kernel is evaluated exactly in the form $(2\pi/i)W_{c} = (uv)^{-1/2}/\big(c(t^{1/2c} - t^{-1/2c})\big)$, $t = u/v$, at points $u = a^{2k}$, $v = b^{2k}$ with $k/c$ an integer ($c = \tfrac13, \tfrac12, 1, 2, 3$): $c = 1$ gives $1/(u - v)$; every kernel and every difference is homogeneous of degree $-1$ under scaling by squares; $D_{c} = 0$ exactly at $c = 1$; and along a ray towards the edge $D_{c}$ grows as the inverse scale; the symbol satisfies the two-point KMS identity at period $2\pi c$ and, for $c \neq 1$, not at $2\pi$; and $t^{1/2c} \mapsto -t^{1/2c}$ under $y \mapsto y - 2\pi i c$ makes the kernel antiperiodic there (the fermionic KMS condition), while at $c = \tfrac12$ the shift by $2\pi i$ leaves it periodic. For $c = \infty$ the symbol is $\chi_{(-\infty, 0)}(w)$ and $(2\pi/i)W_{\infty} = (uv)^{-1/2}/(y + i0)$.

With $S(z) = \sinh z/z$, $g_{c}(y) = y^{-1}\big(S(y/2c)^{-1} - S(y/2)^{-1}\big)$. Both inverses are even power series with constant term $1$, so the pole cancels; $S(z)^{-1} = 1 - \tfrac{z^{2}}{6} + \tfrac{7z^{4}}{360} - \cdots$ gives $g_{c}(y) = -\tfrac{1}{24}(c^{-2} - 1)y + \tfrac{7}{5760}(c^{-4} - 1)y^{3} + \cdots$, and the zeros of $\sinh(y/2c)$ and $\sinh(y/2)$ off $0$ lie at $2\pi i c k$ and $2\pi i k$, $k \neq 0$, so $g_{c}$ is analytic on $|\mathrm{Im}\,y| < 2\pi\min(c, 1)$. The check computes both series exactly to order $12$ in $y$ with $\gamma = c^{-2} \in \{0, \tfrac19, \tfrac14, 1, 4\}$: zero constant term, even, first coefficient $-(\gamma - 1)/24$, identically zero only at $\gamma = 1$; and as a control, omitting the factor $1/c$ leaves the pole $(c - 1)/y$. At a point $u_{0} > 0$, $\lambda W_{0}(u_{0} + \lambda x, u_{0} + \lambda y) = W_{0}(x, y)$ and $\lambda D_{c}(u_{0} + \lambda x, u_{0} + \lambda y) = O(\lambda)$ uniformly for $x, y$ in a compact set, because $D_{c}$ is smooth near $(u_{0}, u_{0})$; for small $\lambda$ both arguments stay on the half $u > 0$, so the cross-edge block does not enter. The scaling limit at $u_{0}$ is $\omega_{0}$, and the same holds at $u_{0} < 0$ by the reflection $u \mapsto -u$. Only the same-half block entered: for test functions supported in a compact set and $\lambda$ small, $u_{0} + \lambda x$ and $u_{0} + \lambda y$ lie on the half containing $u_{0}$. So the conclusion holds for every stationary gauge-invariant quasi-free state with the half restrictions of $\omega_{c}^{\times}$ and an arbitrary cross-edge block (subject only to positivity of the whole symbol), and for $c \neq 1$ such a state is not $\omega_{0}$, because its restriction to $u > 0$ is the I-state at period $2\pi c$ while that of $\omega_{0}$ is the I-state at $2\pi$ (Theorem \claimref{synthesis.positive-null-energy-gives-r-premises}(ii)) and the I-state at a given period is unique (Theorem \claimref{hawking.i-state-continuum}). Across the edge, $\omega_{c}^{\times}(\psi^{*}(v)\psi(u)) = 0$ for $u > 0 > v$ (product of even states, odd factors), while $W_{0}(u, v) = \tfrac{i}{2\pi}(u - v)^{-1} \neq 0$ there, homogeneous of degree $-1$: the difference is unbounded on every neighbourhood of $(0, 0)$, and the check tabulates it at three points and three scales. So $\omega_{c}^{\times} \neq \omega_{0}$ for every $c$; at $c = 1$ both halves are the I-state at period $2\pi$ (the continuum restatement of (R) on each half) and the state is still not the null vacuum.

\emph{(iv)} For a complex chiral fermion $T_{uu} = \tfrac{i}{2}(\psi^{*}\partial_{u}\psi - \partial_{u}\psi^{*}\,\psi)$, renormalised by point-splitting against $\omega_{0}$: $\langle T_{uu}\rangle(u) = \tfrac{i}{2}\lim_{v \to u}(\partial_{u} - \partial_{v})D_{c}(u, v)$. At $u = v = 1$, with $D_{c} = \tfrac{i}{2\pi}X$, the check computes $X(1 + \epsilon, 1) = (1 + \epsilon)^{-1/2}g_{c}(\ln(1 + \epsilon))$ and $X(1, 1 + \epsilon) = (1 + \epsilon)^{-1/2}g_{c}(-\ln(1 + \epsilon))$ as exact series in $\epsilon$ and finds $\langle T_{uu}\rangle(1) = -(X_{u} - X_{v})/(4\pi) = (\gamma - 1)/(48\pi)$ for the five values of $\gamma$ ($-1/(48\pi)$ at $\gamma = 0$, $3/(48\pi)$ at $\gamma = 4$, $0$ at $\gamma = 1$). Since $D_{c}$ is homogeneous of degree $-1$, $\langle T_{uu}\rangle$ is homogeneous of degree $-2$: $(c^{-2} - 1)/(48\pi u^{2})$. The Killing time $s$ along the generator has $du/ds = \kappa u$, so the Killing-frame density is $\kappa^{2}u^{2}\langle T_{uu}\rangle = \kappa^{2}(c^{-2} - 1)/(48\pi)$, constant; at $c = \infty$ it is $-\kappa^{2}/(48\pi)$, the familiar Boulware-type deficit of one complex chiral fermion at $T = \kappa/2\pi$ (cited from memory, as a consistency check only). Symbols enter the kernels linearly, and any measurable symbol with values in $[0, 1]$ defines a gauge-invariant quasi-free state; a function of the Mellin weight alone commutes with the boost, so the state with symbol $\tfrac14 m_{1/2} + \tfrac34 m_{\infty}$ on each half (and the product across the edge) is stationary. Its difference series is the same mix: zero constant term (locally the null vacuum off the edge, as in (iii)), $y$-coefficient of $g$ equal to $-\tfrac{1}{24}\big(\tfrac14(4 - 1) + \tfrac34(0 - 1)\big) = 0$ (zero null energy at every point, by homogeneity), and $y^{3}$-coefficient $\tfrac{7}{5760}\big(\tfrac14(16 - 1) + \tfrac34(0 - 1)\big) = \tfrac{21}{5760} \neq 0$: not $\omega_{0}$. The check verifies these coefficients exactly, that $\tfrac14$ is the only cancelling weight among three tested, and that within the one-period family only $\gamma = 1$ has zero null energy.

\emph{(v)} The first two sentences restate (ii), (iii) and (iv). A condition on the scaling limits at the points $u > 0$, with stationarity, is met by $\omega_{c}^{\times}$ restricted to $u > 0$ for every $c$, which is KMS at $2\pi c$ there, so it implies neither (N$_{+}$) nor (R) on the half-generator, and adding zero null energy still leaves the mixture of (iv). In the store's Painlev\'e-Gullstrand form of the eternal hole the gauge field is invertible on the future horizon $r = 2M$, $u > 0$ (Theorem \claimref{singularity.why-calibration-cannot-be-internal}), while the edge $u = 0$ and the half $u < 0$ lie beyond that chart (Theorem \claimref{synthesis.positive-null-energy-gives-r-premises}, flag beyond one flat chart). On the half-generator the stationary quasi-free states are labelled by their Mellin symbol, a function of $w$, and fixing it to $m_{1}$ is (R) in its continuum form; this is what a principle on the chart that delivers (R) must do. On a local Rindler horizon of Minkowski space the edge point is an ordinary point of the flat chart, and (ii) applies there. Transport-absent: no transport statement is used.

\emph{Remark (not part of the result).} For the outgoing chiral sector of a hole the boost acts on the outgoing affine parameter $U$ as a dilation about $U = 0$, the future horizon itself, which is on the Painlev\'e-Gullstrand chart; the same scaling lemma then concerns a state locally the null vacuum at the future horizon. That case, for a hole formed by collapse, is the main argument of the result filed alongside, synthesis.collapsing-hole-flux-from-ground-state-on-the-ingoing-sheet, and the mechanism of the Fredenhagen-Haag derivation of Hawking radiation (cited from memory). It needs the reduction of the outgoing modes to a two-dimensional chiral field near the horizon, which the store carries only as the postulate J3 (Premise \claimref{hawking.joint-j3-s-wave-reduction}), and asymptotic rather than exact stationarity.
